\documentclass[11pt]{article}

\usepackage[margin=1.1in]{geometry}
\usepackage[utf8]{inputenc}
\usepackage[T1]{fontenc}
\DeclareUnicodeCharacter{03B8}{\ensuremath{\theta}}
\usepackage{amsmath,amssymb}
\usepackage{booktabs}
\usepackage{listings}
\usepackage{xcolor}
\usepackage{microtype}
\usepackage[hidelinks]{hyperref}

% Same plain Julia listing style as queue_layers.tex.
\lstdefinelanguage{Julia}{
  keywords={abstract,type,struct,mutable,function,end,if,else,elseif,for,
            while,begin,return,using,import,export,const,where,in,do,let,
            quote,macro,module,true,false,nothing,missing,local,global},
  sensitive=true,
  morecomment=[l]{\#},
  morestring=[b]",
}
\lstset{
  language=Julia,
  basicstyle=\ttfamily\small,
  keywordstyle=\bfseries,
  commentstyle=\color{black!55}\itshape,
  stringstyle=\color{black!70},
  showstringspaces=false,
  columns=fullflexible,
  keepspaces=true,
  breaklines=true,
  breakatwhitespace=true,
  frame=single,
  framerule=0.3pt,
  xleftmargin=1em,
  xrightmargin=1em,
  aboveskip=0.8\baselineskip,
  belowskip=0.8\baselineskip,
  literate={θ}{{$\theta$}}1
           {τ}{{$\tau$}}1
           {β}{{$\beta$}}1
           {∂}{{$\partial$}}1
           {≤}{{$\le$}}1
           {≥}{{$\ge$}}1
           {∈}{{$\in$}}1
           {∞}{{$\infty$}}1
           {·}{{$\cdot$}}1
           {—}{{---}}1
           {–}{{--}}1,
}

\newcommand{\E}{\mathbb{E}}
\newcommand{\dd}{\mathrm{d}}
\newcommand{\code}[1]{\texttt{#1}}

\title{Concourse: Amendments to the Queueing Model Definition}
\author{Concourse research notes}
\date{July 17, 2026}

\begin{document}
\maketitle
\tableofcontents

%==============================================================================
\section{What this document is}

\code{queue\_layers.tex} (in this directory) defines a three-layer queueing
library: a surface language, a GSMP intermediate representation, and the
CompetingClocks sampler. Concourse.jl adopts that design with the following
changes, and this note is the normative statement of them. Where the two
documents disagree, this one wins; everything not amended here carries over
unchanged.

The amendments, agreed in design discussion on July 17, 2026:

\begin{description}
\item[A1.] \emph{Clock bookkeeping is model state.} Per-clock enabling times
  and banked ages move into \code{QueueState}. The sampler receives them; it
  does not own them. (Section~\ref{sec:a1}.)
\item[A2.] \emph{The clock-family set is open.} \code{queue\_layers.tex}
  claimed two clock families suffice; reneging and server breakdowns refute
  that as a design ceiling. Families become a registry the interpreter
  iterates, not a closed enumeration it hard-codes. (Section~\ref{sec:a2}.)
\item[A3.] \emph{Every function-valued field is a combinator value.} Laws,
  mark samplers, discipline components, and routing kernels are inspectable
  expression values with derivable reads, with one explicitly-marked opaque
  escape hatch. (Section~\ref{sec:a3}.)
\item[A4.] \emph{No hidden \code{rand()}.} Every auxiliary draw --- marks,
  probabilistic routing, SIRO's pick --- consumes an explicit, keyed,
  recorded draw source that appears in the function's signature.
  (Section~\ref{sec:a4}.)
\item[A5.] \emph{State-dependent service laws carry age across segments.}
  A station's service law may read live occupancy; when a watched count
  changes, each surviving clock is re-declared with the fresh law
  conditioned on its accrued service effort, \code{te} fixed. (Added
  July 18, 2026; Section~\ref{sec:a5}.)
\item[A6.] \emph{Round service: the plan is state, the round is one
  clock, the policy is pure.} A station may serve in synchronous rounds
  under a boundary policy; the committed round plan and the per-job work
  counters live in \code{QueueState}, one derived \code{:round} clock per
  station runs each round, and policy purity is a documented modeling
  assumption. (Added July 19, 2026; Section~\ref{sec:a6}.)
\end{description}

One contract delta against ClockGradients.jl falls out of A1 and is proposed
in Section~\ref{sec:cd1}: an optional four-argument
\code{fire(model, state, key, t)}.

A second decision frames the whole note: Concourse depends on
CompetingClocks.jl and ClockGradients.jl and \emph{not} on ChronoSim.jl.
This is cheaper than \code{queue\_layers.tex} assumed, because every seam it
borrowed from ChronoSim exists in the two remaining packages:
\code{:resume} memory is \code{enable!} with a left-shifted enabling time
(\code{te < when} is documented behavior of the sampler interface);
processor sharing's mid-flight re-evaluation is \code{reenable!}; the record
enters the estimators through the \code{gradient\_record} generic, which has
no core methods precisely so that frameworks can attach their own. The one
substantial replacement is the nine-verb branchable-world protocol for the
clone-based estimators, deferred to phase~2. Concourse thereby becomes the
second independent implementation of the ClockGradients model contract,
beside the ChronoSim package extension --- which is the test that the
contract is an abstraction and not a description of its only client.

\paragraph{Falsification stance.} Following the practice in
\code{queue\_layers.tex}~\S5, every design claim in this note appears in
Section~\ref{sec:charter} paired with the concrete phase-1 test that would
falsify it. The phase-1 prototype exists to run that table, not to be a
library.

%==============================================================================
\section{A1: clock bookkeeping is model state}
\label{sec:a1}

\subsection{The argument}

A GSMP's Markov state is $(s, \text{clock readings})$ --- Glasserman Ch.~2
gives the canonical construction. \code{queue\_layers.tex} put $s$ in
\code{QueueState} and left the readings in the sampler. That split fails
twice:

\begin{enumerate}
\item \emph{It breaks \code{fire}'s purity for age-reading disciplines.}
  SRPT's \code{preempt} compares remaining sizes, which are size marks minus
  ages; \code{:resume} re-enabling needs the banked age. If ages live only in
  the sampler, \code{settle!} must call \code{enabled\_ages} on the sampler
  mid-\code{fire}, which inverts the layering and makes \code{fire} depend on
  an argument it does not receive.
\item \emph{It makes the four-argument \code{clock\_distribution} ambiguous
  under replay.} ClockGradients' score replay evaluates
  \code{clock\_distribution(model, θ, k, state)} for every enabled clock $k$
  at every inter-event interval --- at states \emph{later} than $k$'s
  enabling (\code{src/score.jl:53--58}). A time-varying law (nonhomogeneous
  arrivals, a shift schedule) must construct its distribution anchored at
  $k$'s enabling time $t_e$. If the model reads a single ``current time''
  field, a long-enabled clock gets the wrong anchor at later states. Only a
  \emph{per-clock} $t_e$ table inside the state gives the same answer at
  every state from which the replay asks.
\end{enumerate}

The replay machinery already concedes the point: \code{score\_loglikelihood}
maintains an external per-clock enabling-time table via
\code{sync\_enabling\_times!} because the likelihood cannot be computed
without one. A1 moves that table to where the formalism says it lives.

\subsection{The amended state}

\begin{lstlisting}
struct QueueState
    time::Float64                    # time of the firing that produced this state
    buf::Vector{Vector{JobId}}       # per station, in discipline order
    srv::Vector{Vector{JobId}}       # per station, the jobs in service
    jobs::Dict{JobId,Job}            # marks; birth time lives in the record
    pending::Vector{Dict{GroupId,Vector{JobId}}}  # joins: siblings so far
    cursor::Vector{Int}              # stateful routing kernels
    next_id::JobId
    te::Dict{ClockKey,Float64}       # enabling time, per enabled clock
    bank::Dict{ClockKey,Float64}     # accrued age, per disabled :resume clock
end
\end{lstlisting}

Bookkeeping rules, all inside \code{fire}/\code{settle!}:

\begin{itemize}
\item A clock entering the enabled set at firing time $t$ gets
  $\code{te}[k] = t - \code{bank}[k]$ (with $\code{bank}[k] = 0$ when absent),
  and its bank entry is deleted. The left-shifted $t_e$ is exactly what
  \code{enable!(sampler, k, dist, te, when)} receives, so the sampler and the
  state agree by construction.
\item A clock leaving the enabled set without firing, under a discipline with
  \code{memory = :resume}, banks its age:
  $\code{bank}[k] \mathrel{+}= t - \code{te}[k]$. Under \code{:fresh} the
  bank entry is deleted. Firing deletes both entries.
\item A clock's age at any state is $\code{st.time} - \code{st.te}[k]$ for
  enabled clocks and $\code{st.bank}[k]$ for disabled ones. SRPT's
  \code{preempt} and the PS re-enabling read \emph{these}; no model code
  calls \code{enabled\_ages} on the sampler, ever. The sampler is a pure
  consumer of state-derived arguments.
\end{itemize}

Two consequences worth naming. First, \code{states\_equal} and record-fold
replay now compare and reproduce the clock bookkeeping too, so a divergence
between the model's ages and the sampler's is detectable rather than silent
--- the runtime invariant is \code{st.te == sampler's enabling times} after
every firing. Second, the state stays plain data: two \code{Dict}s of
\code{isbits} pairs copy, hash, and diff like everything else in it.

\subsection{Contract delta CD-1: \code{fire} receives the firing time}
\label{sec:cd1}

\code{fire} must stamp \code{st.time} and the new \code{te} entries, so it
needs the firing time --- which the current ClockGradients contract does not
pass: \code{fire(model, state, key)} is called at twelve sites across
\code{score.jl}, \code{spa.jl}, \code{functionals.jl}, \code{clockworld.jl},
\code{records.jl}, and \code{conformance.jl}, and at \emph{every} site the
firing (or swap) time is already in scope in the caller. The proposal:

\begin{lstlisting}
# In ClockGradients: the time-aware form, with a fallback so existing
# time-free models are untouched.
fire(model, state, key, t) = fire(model, state, key)
# All internal call sites switch to the four-argument form.
\end{lstlisting}

Backwards compatible, mechanical, and confined to the user's own package.
Concourse implements only the four-argument method. The alternative ---
keeping \code{fire} time-free and expressing every law in age coordinates ---
was rejected because it silently forbids wall-clock-inhomogeneous laws,
which are a stated goal (``time-varying hazards for rates'').

SPA's hypothetical firings (\code{spa.jl:70,99,271}) pass the swap time;
that is the semantically correct anchor for the alternate world's new
enablings, and the criticality gate's state comparisons inherit it.

%==============================================================================
\section{A2: the clock-family set is open}
\label{sec:a2}

\subsection{The refutation of ``two families suffice''}

\code{queue\_layers.tex} finding~1 claimed arrival and service clocks cover
the design. The claim holds for its chosen feature set and fails one step
outside it: \emph{reneging} --- a patience clock on a job that is waiting,
not in service --- is bread-and-butter queueing (M/M/1+M, every call-center
model) and is a third family. Server breakdowns and vacations are a fourth.
The finding is therefore demoted from an architectural assumption to a
statement about the initial feature set, and the architecture must make
adding a family cheap.

\subsection{The family registry}

\code{ClockKey} stays one concrete isbits tuple; the \code{Symbol} field is
the open dimension:

\begin{lstlisting}
const ClockKey = Tuple{Symbol,Int32,Int64}
# (:arrival,  station, 0)      the source's next-arrival clock
# (:service,  station, jobid)  one clock per job in service
# (:patience, station, jobid)  one clock per waiting job, if reneging   [later]
# (:repair,   station, srvidx) one clock per down server, if breakdowns [later]
\end{lstlisting}

A clock family is a value with three functions --- the same move A3 makes
for laws, applied to the interpreter itself:

\begin{lstlisting}
struct ClockFamily
    name::Symbol
    enumerate  # (model, state) -> iterator of ClockKey   (who is enabled)
    law        # (model, θ, key, state) -> UnivariateDistribution
    fire       # (model, state, key, t) -> new state      (what firing means)
end
\end{lstlisting}

\code{QueueGSMP} carries \code{families::Vector{ClockFamily}};
\code{enabled} is the concatenation of each family's \code{enumerate} in
registry order (deterministic order, as the contract requires);
\code{clock\_distribution} and \code{fire} dispatch on \code{key[1]} through
the registry. The interpreter, the record fold, the checker's randomness
census, and the estimators never mention a family by name. The acid test is
F2 in the charter: adding \code{:patience} must touch the registry and
nothing else.

What stays closed: a family's \code{fire} still works through the shared
\code{deposit!}/\code{settle!} vocabulary, so disciplines and routing
compose with new families instead of each family reinventing station
semantics. (A patience firing is ``remove job $j$ from \code{buf}, deposit
to the reneging destination''; a repair firing is ``increment station
capacity, \code{settle!}''.)

%==============================================================================
\section{A3: every function-valued field is a combinator value}
\label{sec:a3}

\subsection{The argument}

The IR's advertised virtue is that it is \emph{data} --- printable,
diffable, hashable, checkable. \code{CompiledStation}'s
\code{service::Function}, \code{routing::Function}, \code{mark::Function}
fields quietly forfeit that: closures cannot be hashed, compared,
serialized, or introspected, and the checks that justify the IR ---
C2 (mark typing) and C6 (randomness census / estimator tiers) --- need to
know what a law reads and whether it depends on $\theta$. The routing
kernels already show the fix (\code{Always}, \code{ByMark},
\code{JoinShortestQueue} are values); A3 applies it uniformly.

\subsection{The scalar-expression algebra}

One small expression type feeds every parameter slot:

\begin{lstlisting}
# ScalarExpr: a tree with derivable reads.
Param(:mu)          # reads θ by declared name -> θ-dependence is visible
Mark(:size)         # reads a job mark        -> C2 dataflow is visible
Enab()              # the wall-clock enabling time t_e
Const(2.0)
+, *, inv, log, ... # lifted arithmetic on ScalarExpr
\end{lstlisting}

A \emph{law} is a distribution family applied to scalar expressions; a
\emph{mark sampler} is a named tuple of distribution expressions; a
\emph{discipline} is the four-field value of \code{queue\_layers.tex}
\S2.2 with its \code{by}-style components as \code{ScalarExpr} rather than
closures:

\begin{lstlisting}
service = Law(Exponential, scale = Mark(:size) * inv(Param(:mu_short)))
mark    = MarkLaw(size = Law(LogNormal, μ = Const(0.0), σ = Const(1.0)))
disc    = Priority(by = Mark(:class); preempt = true, memory = :resume)
\end{lstlisting}

Derived, not declared: \code{reads\_θ(law)}, \code{reads\_marks(law)},
\code{reads\_time(law)} walk the tree. C2 becomes a dataflow check over
values; C6's census and the estimator-tier table are computed, not asserted.
Nonhomogeneous laws use \code{Enab()}; state-reading kernels
(\code{JoinShortestQueue}) remain the named combinators they already were.

\subsection{The escape hatch}

\begin{lstlisting}
OpaqueLaw(f; reads_θ = true, reads_marks = (:size,), reads_time = false)
\end{lstlisting}

An opaque law carries \emph{mandatory} declared reads; the checker trusts
the declaration, marks every conclusion resting on it as
``declared, unverified,'' and degrades the estimator tiers it cannot
certify. A closure with no declaration is a construction-time error, not a
warning. The same wrapper exists for kernels and discipline components. The
charter's F3 measures whether the algebra is rich enough that the hatch
stays shut for the standard models.

%==============================================================================
\section{A4: no hidden randomness}
\label{sec:a4}

The rule from \code{queue\_layers.tex}~\S3.5 --- every draw is a clock, a
recorded mark, or a recorded choice --- becomes a signature-level fact.
Auxiliary draws (marks at arrival, \code{Probabilistic} routing, SIRO's
pick) consume an explicit draw source:

\begin{lstlisting}
# The draw source is an argument wherever auxiliary randomness may occur.
select(disc, buf, state, draws)          # SIRO reads; FCFS ignores
route(kernel, job, state, draws)         # Probabilistic reads; ByMark ignores
draw_marks(marklaw, draws)               # at arrival firings
\end{lstlisting}

Two modes, one type. In live simulation, \code{draws} wraps
CompetingClocks' \code{KeyedStreams}, keyed by
\code{(firing ClockKey, purpose::Symbol)} so that a draw belongs to a clock
and purpose rather than to a position in a global call sequence --- the
same common-random-numbers discipline the sampler already imposes on clock
randomness, extended to the auxiliary draws. Every value drawn is appended
to the current record entry. In replay, \code{draws} is the recorded
sequence, and drawing returns the recorded values in consumption order ---
which is precisely the conditioning that makes \code{fire} deterministic
given the record, and which the likelihood's mark and routing-mass terms
require.

The marked record is therefore
\[
\mathcal{R} = \big( (k_i, t_i, a_i) \big)_{i},
\qquad a_i = \text{the auxiliary draws consumed by firing } i,
\text{ in order},
\]
\code{queue\_layers.tex}'s record with the mark column generalized to
``all auxiliary draws.'' The checker's census (C6) enumerates the possible
draws per family from the combinator values (A3), so
record-schema-vs-model agreement is itself checkable.

%==============================================================================
\section{A5: state-dependent service laws carry age across segments}
\label{sec:a5}

Added July 18, 2026, with the state-dependent-service capability; the
phase-1 amendments above are unchanged by it.

A station's service law may read live occupancy through two new atoms of
the A3 algebra, \code{InService(station)} and \code{InBuffer(station)},
with \code{reads\_state} joining the derived censuses. Occupancy is live
state, not a value frozen at enabling, so the read is scoped: only the
service law of a \emph{station} may read occupancy. State in a mark law
would push station state into the record's mark draws, violating A4's
conditioning; a state-reading interarrival law is balking, a feature of
its own; patience laws, routing kernels, and discipline ordering keys are
refused likewise. This is the compile check the code registry numbers C5
(see the registry remark below).

The semantics is the \emph{age-carryover segment convention}. Compile
builds a dependency map (\code{srv\_readers}/\code{buf\_readers}) from the
census. When a firing changes a watched in-service or buffer count at
time $t$, every surviving in-service clock of a reader station closes a
segment: it banks the effort accrued since its last anchor at the old
speed ($\min(1, \code{servers}/n)$ under PS, $1$ elsewhere), moves its
anchor to $t$, keeps \code{te} at the original enabling --- \code{te} is
never rewritten --- and is re-declared to the sampler as
\code{(:reenable, k)}. The re-declared law is the \emph{fresh base}
$F_{\mathrm{new}}$, rebuilt from the occupancy at $t$, conditioned on the
total draw exceeding the accrued effort $a$: $X \sim F_{\mathrm{new}}
\mid X > a$, delivered as a wall-time distribution anchored at $t$ (the
existing \code{SharedRemaining} wrapper). Processor sharing becomes a
special case of this general path: a PS station is an srv-reader of
itself, with base unchanged and speed $\min(1, \code{servers}/n)$.

The rejected alternative was the quantile-equivalent (hazard-matching)
mapping, which carries the old law's survival quantile into the new law
so that survival probability is continuous across the change. Age
carryover won on two grounds: it is what processor sharing already
implements, and it is what the ClockGradients fixed-anchor contract
expects --- \code{te} anchors every likelihood segment, and the score
replay rebuilds the current segment's law at every state it visits.
Pathwise branching is \emph{not} admitted for these models
(\code{branch\_world} refuses): reordering events changes an occupancy,
hence the law, discontinuously, and the clone estimators' unbiasedness
argument does not cover that coupling. The score estimator remains valid;
the stability check (C4) skips state-reading stations rather than invent
a static mean.

\paragraph{Registry remark.} \code{queue\_layers.tex}~\S3.7 numbers its
checks C1--C6 with C5 = oracle detection and C6 = randomness census, and
this note cites those meanings above (Sections~\ref{sec:a3},
\ref{sec:a4}, \ref{sec:charter}). Concourse's \emph{compile-time} registry
diverged: its C5 is the state-reading scope check of this amendment (the
oracle detector is not a compile check in Concourse). In this note,
``compile check C5 (state-reading scope)'' refers to the code's registry;
an unqualified C5 elsewhere in this document keeps the
\code{queue\_layers.tex} meaning.

%==============================================================================
\section{A6: synchronous round service}
\label{sec:a6}

Added July 19, 2026, with the round-based token service capability
(iteration-level batching for LLM serving; Dai et al.\ arXiv:2504.07347
and Dong \& Cao arXiv:2604.11001 are the target models). A1--A5 are
unchanged by it. A round station serves in synchronous rounds: at each
boundary a policy examines the waiting line and the active set, admits and
evicts jobs, and allocates integer work units to the active ones; one
station-level clock runs the round; at its firing the allocations are
credited against per-job work counters, exhausted jobs route onward, and
the next round is planned at the same instant. Three commitments define
the amendment.

\paragraph{The round plan is model state.} A1 applied again:
\code{QueueState} gains \code{work} (per-job remaining integer work, one
counter per declared phase, created at admission and deleted at
departure) and \code{roundplan} (one \code{Union\{RoundPlan,Nothing\}}
slot per station holding the running round's frozen admissions,
evictions, per-job allocations, and aggregate pseudo-marks). Both are
replay-owned exactly like \code{cursor}, \code{te}, and \code{bank}:
\code{states\_equal} compares them field by field, and the record fold
reproduces them, so the round machinery adds nothing to the record beyond
the draws a policy explicitly consumes (A4).

\paragraph{One clock per round, membership derived.} The \code{:round}
family joins the A2 registry with key
\code{(:round, station, 0)} and derived membership: the clock is enabled
exactly while the station's \code{roundplan} slot is non-empty, so the
derived-delta pass --- not the planner --- owns the clock lifecycle. The
round clock \emph{is} the station's single server: active jobs hold no
per-job clocks, the duration law reads only the committed plan's
aggregate pseudo-marks (frozen at enabling, like every mark), and the
firing credits the plan, routes the finished jobs, clears the slot, and
re-plans immediately at the same time, so back-to-back rounds chain with
no gap. A station with no plan idles clock-free, and the first deposit
after an idle period re-anchors the round grid at that deposit's time ---
a documented modeling convention, invisible in the papers' loaded
regimes.

\paragraph{Policy purity is a modeling assumption.} The policy hook
\code{plan\_round(policy, view, draws)} is the first user code inside the
fire path. Its contract is A4 extended to user code: the policy is a pure
function of \code{(view, draws)} --- all randomness through the recorded
draw source, no clocks, no wall time, no $\theta$, no state beyond the
read-only view. Like an \code{Opaque} law's declared reads (A3), purity
cannot be compile-checked for arbitrary Julia code; it is an assumption
the model owner takes on. The view is constructed as a barrier (copies
and values, never \code{QueueState}) so that honest code cannot violate
the contract by accident, and replay equality (F4) plus the debug
membership oracle are the enforcement that exists --- the test suite
demonstrates both catching a deliberately impure policy.

%==============================================================================
\section{The amended contract surface}
\label{sec:contract}

What Concourse implements, in ClockGradients' vocabulary. Deltas against the
shipped contract are marked.

\begin{center}
\small
\begin{tabular}{@{}lll@{}}
\toprule
Contract function & Concourse implementation & Delta \\
\midrule
\code{initial\_state(m)} & empty places; \code{te}, \code{bank} empty; \code{time = 0} & --- \\
\code{clockkeytype(m)} & \code{Tuple\{Symbol,Int32,Int64\}} & --- \\
\code{enabled(m, st)} & concat of family \code{enumerate}s, registry order & --- \\
\code{clock\_distribution(m, θ, k, st)} & family \code{law}; anchors from \code{st.te} & --- \\
\code{fire(m, st, k, t)} & family \code{fire} through \code{deposit!}/\code{settle!} & CD-1 \\
\code{gradient\_record(m, rec, θ)} & one method for the marked record type & by design \\
\bottomrule
\end{tabular}
\end{center}

Replay determinism is conditional on the recorded draws (A4), exactly as
the ChronoSim extension conditions on its recorded firings. The nine-verb
branchable world over \code{(QueueGSMP, QueueState, SamplingContext)} is
phase~2; \code{check\_branchable} is its acceptance test.

%==============================================================================
\section{The falsification charter}
\label{sec:charter}

Each claim is stated so that the phase-1 (or explicitly phase-2) prototype
can refute it. A failed test is a design discovery, not a bug to paper over.
Oracles follow the 4-standard-error convention; closed forms come from the
C5 detector cases (M/M/1, the SITA farm as two independent M/G/1 queues,
Erlang-A for reneging).

\begin{description}
\item[F1 (A1 sufficiency).] With \code{te}/\code{bank} in state, no model
  code needs sampler introspection. \emph{Test:} SRPT and
  preempt-\code{:resume} runs match their oracles while the model layer has
  no reference to \code{enabled\_ages}; the \code{st.te}-vs-sampler
  invariant holds after every firing of a long mixed run.
\item[F2 (A2 openness).] A new clock family costs one registry entry.
  \emph{Test:} implement \code{:patience} (reneging); diff the change ---
  if any file other than the family registry and the surface constructor
  changes, the claim is falsified. M/M/1+M against Erlang-A.
\item[F3 (A3 coverage).] The scalar algebra covers the five
  \code{queue\_layers.tex} models plus reneging with zero \code{Opaque}
  escapes. \emph{Test:} build all six; count escapes.
\item[F4 (A4 determinism).] \code{fire} is deterministic given the record.
  \emph{Test:} fold the marked record and compare \code{states\_equal}
  against the live trajectory at every index, including \code{te} and
  \code{bank}.
\item[F5 (CD-1 sufficiency).] The four-argument \code{fire} is the only
  ClockGradients change queueing needs. \emph{Test:} score and IPA on a
  nonhomogeneous-arrival M(t)/M/1 against finite differences, running on
  patched ClockGradients with no other core edits.
\item[F6 (cascade totality).] Under check C3 (acyclic blocking),
  \code{deposit!}/\code{settle!} cascades terminate deterministically.
  \emph{Test:} property test over random finite-buffer networks with
  \code{:block}; assert termination and replay-order equality. (The
  worklist order itself is a middle-layer decision,
  Section~\ref{sec:queued}.)
\item[F7 (initial family set).] \code{\{:arrival, :service\}} suffices for
  the five \code{queue\_layers.tex} models --- the surviving, weakened form
  of the old finding~1. \emph{Test:} phase 1 builds all five with two
  families; reneging (F2) is the designed refutation lying one step beyond.
\item[F8 (truncate-and-rescale, phase 2).] The speed-free compilation of PS
  equals an explicit-speed reference. \emph{Test:} M/M/1-PS against its
  closed form, and against a Bogliolo-style clock-and-speed-matrix
  reference implementation on a small case.
\end{description}

\paragraph{Compile-check registry (continued).} The registry remark of
Section~\ref{sec:a5} records that Concourse's compile-time check numbering
diverged from \code{queue\_layers.tex} at C5. The July 2026 capabilities
extend the code's registry; this charter is the normative record of the
numbering:

\begin{description}
\item[C11 (remark scope).] Deposit-time mark-redraw (\code{remark}) laws
  obey source-mark-law scope: no station state, and reads only of the
  marks the job carries \emph{before} the redraw. Their draws ride the
  depositing firing's draw list (A4); a remark-only mark is readable at
  the remark station and downstream, never upstream.
\item[C12 (round station shape).] A round station composes with exactly
  the plain FCFS shape: non-preemptive \code{:back}-insert FCFS, no
  \code{batching} (the policy owns batch formation), \code{servers = 1}
  (the round clock is the server, A6), no \code{service} law (the
  duration lives in the \code{Rounds} config), and \code{work} marks
  produced upstream --- integer-valued and nonnegative at admission,
  checked at runtime.
\item[C13 (round duration-law scope).] The round duration law reads only
  parameters, the enabling time, and the committed plan's aggregate
  pseudo-marks --- never job marks and never live state.
\end{description}

Every message is pinned verbatim by the test suite
(\code{test/test\_remark.jl}, \code{test/test\_rounds.jl}). The
\code{ShortestQueue} destination checks (destinations must exist; under
\code{by = :tokens} every destination must be a round station) are plain
\code{ArgumentError}s and deliberately carry no C~number.

%==============================================================================
\section{Queued middle-layer questions}
\label{sec:queued}

Decisions that belong to the event-loop design conversation, deliberately
not made here. A1/CD-1 already settled ``who owns time'' (the state records
it; the interpreter advances it; the sampler is told it).
\emph{Update, later the same day:} items 1, 3, and 4 are decided and item 2
is framed with a proposal in \code{event\_loop.tex} (this directory), which
is now the normative middle-layer document.

\begin{enumerate}
\item \emph{Diff vs.\ push clock lifecycle.} Recompute \code{enabled} and
  diff after each firing (simple, $O(\text{clocks})$ per event), or have
  \code{fire} emit its enable/disable deltas (\code{fire\_changes} exists in
  the contract for this)? Phase 1 should measure before optimizing.
\item \emph{The cascade worklist.} Order and representation for multi-hop
  blocking chains; the determinism argument F6 must be made for whatever is
  chosen.
\item \emph{Record capture.} Interpreter-owned marked record vs.\ a
  CompetingClocks watcher middleware with an auxiliary-draw column bolted
  on.
\item \emph{Sampler spec policy.} \code{NextReactionMethod} default,
  \code{FirstToFireMethod(coupling=:carry)} when IPA wants the carry
  coupling --- who chooses, and where the choice surfaces in the API.
\end{enumerate}

\end{document}
